\documentclass[10pt,a4paper]{amsart}
\usepackage[margin=19mm]{geometry}
\usepackage{amsmath,amssymb,mathtools}
\usepackage[scr=rsfs,cal=euler]{mathalfa}
\usepackage{xfrac}
\usepackage[unicode,hidelinks,bookmarks=false]{hyperref}
\newcommand{\ie}{i.e.\ }
\newcommand{\cf}{cf.\ }
\newcommand{\resp}{respectively\ }
\newcommand{\Subsection}{\subsection}
\providecommand{\nobreakdash}{\nobreak\hbox{-}}
\setlength{\emergencystretch}{2em}
\multlinegap0pt
\usepackage{bm}
\usepackage{bbm}
\usepackage{dsfont}
\usepackage{xspace}
\usepackage{cancel}
\usepackage{mathtools}
\usepackage{amsthm}
\makeatletter
\@ifundefined{theorem}{\newtheorem{theorem}{Theorem}}{}
\@ifundefined{claim}{\newtheorem{claim}[theorem]{Claim}}{}
\makeatother
\title[On the dib-chromatic number of a digraph]{On the dib-chromatic number of a digraph}
\author{Juan Jos\'e Montellano-Ballesteros, Christian Rubio-Montiel}
\date{Source: arXiv:2511.08807v1 (CC BY 4.0)}
\begin{document}
\maketitle

\noindent\emph{Source and licence.} On the dib-chromatic number of a digraph. Version arXiv:2511.08807v1 (CC BY 4.0), distributed under the Creative Commons Attribution 4.0 International licence, as recorded in the arXiv licence metadata for this exact version and shown on the abstract page https://arxiv.org/abs/2511.08807v1. Full rights notice: licenses/arxiv-2511.08807-CC-BY-4.0.txt.\par\smallskip

\noindent\textit{Editorial scope.} Counted unit: the Theorem whose source label is \texttt{d=3} in the article (source0.tex lines 218--219, source label \texttt{d=3}), delivered together with the article's own complete original proof of it (source0.tex lines 220--291, one \texttt{proof} environment of 72 lines). No mathematics is written, repaired or re-derived: statement and proof are the article's own text line for line. Required context retained uncounted: none. Every definition, display and label of those lines is retained unchanged. Omitted from this package and not redistributed: 1--217, 292--429. Nothing inside a retained excerpt is omitted or altered: every retained body is delivered exactly as the article's lines are written, so no omission receipt of any kind accompanies this package. Citations: the retained excerpts carry no citation command at all, so no bibliography entry of the article is transcribed and no reference is redistributed. Reference adaptation: none was needed and none was made; every reference inside the delivered unit resolves inside it, so no unresolved reference or citation is produced. Printed slips kept literal: no printed word, symbol, hypothesis or number of the counted statement or of its proof was altered. Layout and class: the article's own class is \texttt{article} with packages including amssymb, amsmath, amsthm, fullpage, graphicx, xcolor, color; the wrapper is the lane's pinned \texttt{amsart} editorial wrapper at 10pt on A4, which restates the theorem-like environments the retained text needs and adds the article's own macro definitions that the retained text needs (none), with extra wrapper packages (bm, bbm, dsfont, xspace, cancel, mathtools). The delivered numbering is the unit's own. Assets: no retained span contains a figure environment, a graphics-inclusion command, a PDF-page inclusion, a table or a TikZ picture; no image file is copied into this package, the package declares no asset dependency and no image file is redistributed with it. Licence: version arXiv:2511.08807v1 of this article is distributed under the Creative Commons Attribution 4.0 International licence, as recorded in the arXiv licence metadata for this exact version and shown on the abstract page https://arxiv.org/abs/2511.08807v1. A full rights notice is delivered at licenses/arxiv-2511.08807-CC-BY-4.0.txt. Characters: every retained excerpt is pure ASCII, so the delivered engine, XeLaTeX with its default Latin Modern text font, reports no missing character.
\begin{theorem}\label{d=3} Let $D$ be a simple bipartite  digraph.  If $\delta(D)\geq 2$ then $\text{dib}(D)\geq 3$.     
\end{theorem} 

\begin{proof} Let $D=(A,B)$ be a simple bipartite  digraph, with $\delta(D)\geq 2$ and suppose  $\text{dib}(D)\leq 2$. 
A (non-directed) path $P = (x_1, \dots, x_k)$ in $D$ will be called \emph{bad} if every $x_{2i}\in V(P)$, with $1\leq 2i\leq k$, $x_{2i}$ is a sink in $P$ (respectively, $x_{2i}$ is a source in $P$); and every $x_{2i-1}\in V(P)$, with $1\leq 2i-1\leq k$, $x_{2i-1}$ is a source in $P$ (respectively, $x_{2i-1}$ is a sink in $P$).




Let $P=(x_1, \dots, x_k)$ be a bad path in $D$ of maximum order. Without loss of generality, we will assume that every $x_{2i}\in V(P)$, with $1\leq 2i\leq k$, $x_{2i}$ is a sink in $P$; every $x_{2i-1}\in V(P)$, with $1\leq 2i-1\leq k$, $x_{2i-1}$ is a source in $P$ and that  the sources are in $A$ and the sinks are in $B$.


\begin{claim}\label{b8} $k\leq 7$.
\end{claim}
\begin{proof} Suppose there is a bad path $(x_1, x_2, \dots, x_k)$ in $D$ with $k\geq 8$.   Let $\Gamma: V(D) \rightarrow \{c_1, c_2, c_3\}$ defined as: $\Gamma(x_1) = \Gamma(x_4) = \Gamma(x_7) = c_1$; $\Gamma(x_2) = \Gamma(x_5) = \Gamma(x_8)= c_2$ and $\Gamma(x_3) = \Gamma(x_6) = c_3$. Also, for every $x\in A\setminus\{x_1, x_3, x_5, x_7\}$, $\Gamma(x) = c_1$ and for every $x\in B\setminus\{x_2, x_4, x_6, x_8\}$, $\Gamma(x) = c_3$. It is not hard to see that this is a $b$-coloring of $D$ with 3 colors, where $x_2, x_4$ and $x_6$ are $b^-$-vertices and $x_3, x_5$ and $x_7$ are $b^+$-vertices,  which is not possible. 
\end{proof}


\begin{claim}\label{b7} $k\leq 6$.
\end{claim}
\begin{proof} Suppose there is a bad path $P=(x_1, x_2, \dots, x_7)$ in $D$. By Claim \ref{b8} there is no bad path of order $8$, thus $N^+(x_7) \subseteq V(P)$ and since $\delta(D) \geq 2$,  either $x_7x_2\in E(D)$ or $x_7x_4\in E(D)$. If $x_7x_2\in E(D)$ then in $D$ there is the undirected cycle $C_6= (x_2, x_3, x_4, x_5, x_6, x_7)$.  Let $\Gamma: V(D) \rightarrow \{c_1, c_2, c_3\}$ defined as: $\Gamma(x_3) = \Gamma(x_6) = c_1$; $\Gamma(x_4) = \Gamma(x_7) = c_2$ and $\Gamma(x_2) = \Gamma(x_5) = c_3$. Also, for every $x\in A\setminus\{x_3, x_5, x_7\}$, $\Gamma(x) = c_1$ and for every $x\in B\setminus\{x_2, x_4, x_6\}$, $\Gamma(x) = c_3$. It is not hard to see that this is a $b$-coloring of $D$ with 3 colors, where $x_2, x_4$ and $x_6$ are $b^-$-vertices and $x_3, x_5$ and $x_7$ are $b^+$-vertices, which is not possible. 

Hence, $x_7x_4\in E(D)$ and, by an analogous argument than above, $x_1x_4\in E(D)$. Thus, $\{x_1, x_3, x_5, x_7\}\subseteq N^-(x_4)$ and since $\delta(D) \geq 2$, there is a pair $\{z_1, z_2\}\subseteq N^+(x_4)$ in $A\setminus\{x_1, x_3, x_5, x_7\}$. 
Let $\Gamma: V(D) \rightarrow \{c_1, c_2, c_3\}$ defined as: $\Gamma(x_1) = \Gamma(x_4)= \Gamma(x_5)= c_1$; $\Gamma(x_2) = \Gamma(x_7) =\Gamma(z_1) = c_2$ and $\Gamma(x_3) = \Gamma(x_6) = \Gamma(z_2) = c_3$. Also, for every $x\in A\setminus\{z_1, z_2, x_1, x_3, x_5, x_7\}$, $\Gamma(x) = c_3$ and for every $x\in B\setminus\{x_2, x_4, x_6\}$, $\Gamma(x) = c_2$. It is not hard to see that this is a $b$-coloring of $D$ with 3 colors, where $x_4$ is a $b$-vertex, $x_2$ and  $x_6$ are $b^-$-vertices; and $x_3$ and $x_7$ are $b^+$-vertices, which is not possible. 
\end{proof}


\begin{claim}\label{b6} $k\leq 5$.
\end{claim}
\begin{proof} Suppose there is a bad path $P=(x_1, x_2, \dots, x_6)$ in $D$. By Claim \ref{b7} there is no bad path of order $7$, thus $N^+(x_1) \subseteq V(P)$ and since $\delta(D) \geq 2$,  either $x_1x_6\in E(D)$ or $x_1x_4\in E(D)$. If $x_1x_6\in E(D)$ then in $D$ there is the undirected cycle $C_6= (x_1, x_2, x_3, x_4, x_5, x_6)$. As in Claim \ref{b7} we see that this is not possible. Thus  $x_1x_4\in E(D)$. In an analogous way we see that $x_3x_6\in E(D)$.  Let $\Gamma: V(D) \rightarrow \{c_1, c_2, c_3\}$ defined as: $\Gamma(x_1) = \Gamma(x_6) = c_1$; $\Gamma(x_3) = \Gamma(x_4) = c_2$ and $\Gamma(x_2) = \Gamma(x_5) = c_3$. Also, for every $x\in A\setminus\{x_1, x_3, x_5\}$, $\Gamma(x) = c_1$ and for every $x\in B\setminus\{x_2, x_4, x_6\}$, $\Gamma(x) = c_3$. It is not hard to see that this is a $b$-coloring of $D$ with 3 colors, where $x_1, x_3$ and $x_5$ are $b^+$-vertices; and $x_2, x_4$ and $x_6$ are $b^-$-vertices, which is not possible. 
\end{proof}

\begin{claim}\label{b5} $k\leq 4$.
\end{claim}
\begin{proof} Suppose there is a bad path $P=(x_1, x_2, \dots, x_5)$ in $D$. By Claim \ref{b6} there is no bad path of order $6$, thus $N^+(x_1) \subseteq V(P)$ and since $\delta(D) \geq 2$,  $x_1x_4\in E(D)$. In an analogous way, we see that $x_5x_2\in E(D)$. Moreover, $d^+(x_1) = d^+(x_3) = d^+(x_5)=2$. Thus, $\{x_1, x_3, x_5\}\subseteq N^-(x_2)$ and $\{x_1, x_3, x_5\}\subseteq N^-(x_4)$ and therefore $N^+(x_2), N^+(x_4)\subseteq A\setminus \{x_1, x_3, x_5\}$. Let $\{z_1, z_2\}\subseteq N^+(x_2)$ and $\{w_1, w_2\}\subseteq N^+(x_4)$. 

\noindent {\bf Case 1.} $z_1= w_1$ and $z_2 = w_2$.

\noindent Since $\delta(D) \geq 2$ and $\{x_2, x_4\}\subseteq N^-(z_1)$, there is a pair $\{y_1, y_2\}\subseteq N^+(z_1)$ such that $\{y_1, y_2\} \subseteq B\setminus\{x_2, x_4\}$. Let $\Gamma: V(D) \rightarrow \{c_1, c_2, c_3\}$ defined as: $\Gamma(x_2) =\Gamma(x_5) = \Gamma(z_1) = c_1$; $\Gamma(x_1)  = \Gamma(x_4)  = \Gamma(y_2) =  c_2$ and $\Gamma(x_3) = \Gamma(y_1) = \Gamma(z_2) = c_3$. Also, for every $x\in A\setminus\{z_1, z_2, x_1, x_3, x_5\}$, $\Gamma(x) = c_1$ and for every $x\in B\setminus\{x_2, x_4, y_1, y_2\}$, $\Gamma(x) = c_2$. It is not hard to see that this is a $b$-coloring of $D$ with 3 colors, where $x_4$ is a  $b$-vertex,  $x_3$ and $z_1$ are $b^+$-vertices; and $x_2$ and $z_2$ are $b^-$-vertices,  which is not possible. 

\noindent {\bf Case 2.} $z_1= w_1$ and $z_2 \not=  w_2$.

\noindent  Let $\Gamma: V(D) \rightarrow \{c_1, c_2, c_3\}$ defined as: $\Gamma(x_1) = \Gamma(x_4) = \Gamma(z_2) = c_1$; $\Gamma(x_5)  = \Gamma(z_1) = c_2$ and $\Gamma(x_2) = \Gamma(x_3) = \Gamma(w_2) = c_3$. Also, for every $x\in A\setminus\{z_1, z_2, w_2, x_1, x_3, x_5\}$, $\Gamma(x) = c_3$ and for every $x\in B\setminus\{x_2, x_4\}$, $\Gamma(x) = c_1$. It is not hard to see that this is a $b$-coloring of $D$ with 3 colors, where $x_2$ and $x_4$ are  $b$-vertex, and $x_5$ and $z_1$ is a $b^+$-vertex and a $b^-$-vertex, respectively; which is not possible. 

\noindent {\bf Case 3.}  $\{z_1, z_2\}\cap \{w_1, w_2\} = \emptyset$.

\noindent  Since $w_1 \not \in N^+(x_2)$ and  $\delta(D) \geq 2$ there is $y_1\in N^-(w_1)\cap (B\setminus \{x_2, x_4\})$. 
Let $\Gamma: V(D) \rightarrow \{c_1, c_2, c_3\}$ defined as: $\Gamma(x_1) =\Gamma(x_4) = \Gamma(z_1) = c_1$; $\Gamma(x_2)  = \Gamma(x_3)  = \Gamma(y_1) =\Gamma(w_2) =  c_2$ and $\Gamma(x_5) = \Gamma(w_1) = \Gamma(z_2) = c_3$. Also, for every $x\in A\setminus\{z_1, z_2, w_1, w_2, x_1, x_3, x_5\}$, $\Gamma(x) = c_1$ and for every $x\in B\setminus\{x_2, x_4, y_1\}$, $\Gamma(x) = c_2$. It is not hard to see that the chromatic classes of colors $c_1$ and $c_3$ are acyclic. For the chromatic class of color $c_2$, observe that the only vertices of color $c_2$ in $A$ are $x_3$ and $w_2$. Thus, any directed cycle of color $c_2$ is a square and contains the vertices $\{x_3, w_2\}$. Since $d^+(x_3)= 2$ and $\Gamma (x_4)= c_1$, any directed cycle of color $c_2$ must contain the path $(x_3, x_2, w_2)$ which is not possible since $x_2w_2\not\in E(D)$. Hence, $\Gamma$ is a $b$-coloring of $D$ with 3 colors, where $x_2$ and $x_4$ are  $b$-vertices, and $x_5$ and $w_1$ is a $b^+$-vertex and a $b^-$-vertex, respectively; which is not possible.
\end{proof}

\begin{claim}\label{2reg} $D$ is $2$-regular.
\end{claim}
\begin{proof} Let $x\in V(D)$ with $d^+(x) =r\geq 2$ and let $N^+(x) = \{y_1, \dots, y_r\}$. By Claim \ref{b5} it follows that for some $z\in V(D)\setminus \{x, y_1, \dots, y_r\}$,  $\bigcap\limits_{y\in N^+(x)} N^-(y) = \{x, z\}$, which again, by Claim \ref{b5}, implies that $r=2$ (in other case, there is a bad path of order $5$ where $x_1, x_3$ and $x_5$ are sinks, and $x_2$, $x_4$ are sources). Thus, for every $x\in V(D)$, $d^+(x) =2$ and since $\delta(D) \geq 2$, the claim follows.
\end{proof}

\begin{claim}\label{b4} There are no bad paths of order $4$.
\end{claim}
\begin{proof}
Let $P=(x_1, x_2, x_3, x_4)$ be a bad path in $D$. By Claim \ref{b5} it follows that $x_1x_4\in E(D)$. By Claim \ref{2reg}, $\delta(D) = 2$ and therefore we see that $N^+(x_2) = \{z_1, z_2\}$ and $N^+(x_4) = \{w_1, w_2\}$ where $N^+(x_2), N^+(x_4) \subseteq A\setminus\{x_1, x_3\}$. Again, by Claim \ref{b5}, either $w_1= z_1$ and $z_2=w_2$;  or $\{z_1, z_2\}\cap \{w_1, w_2\}=\emptyset$. 

\noindent {\bf Case 1.} $\{z_1, z_2\}\cap \{w_1, w_2\}=\emptyset$.

\noindent By Claim \ref{b5} there is a pair  $\{y_1, y_2\}\subseteq B\setminus \{x_2, x_4\}$ such that $N^+(y_1) = \{z_1, z_2\}$ and $N^+(y_2) = \{w_1, w_2\}$. Let $\Gamma: V(D) \rightarrow \{c_1, c_2, c_3\}$ defined as: $\Gamma(x_1) =\Gamma(y_2) = \Gamma(z_1) = \Gamma(w_1) = c_1$; $\Gamma(x_2)  = \Gamma(x_3)  = \Gamma(w_2) = c_2$ and $\Gamma(x_4) = \Gamma(y_1) = \Gamma(z_2) = c_3$.  Also, for every $x\in A\setminus\{z_1, z_2, w_1, w_2, x_1, x_3, x_5\}$, $\Gamma(x) = c_2$ and for every $x\in B\setminus\{x_2, x_4, y_1, y_2\}$, $\Gamma(x) = c_3$. It is not hard to see that $\Gamma$ is a $b$-coloring of $D$ with 3 colors, where  $x_4$ is a  $b$-vertex;  $x_1$ and $x_2$ are $b^+$-vertices; and $z_1$ and $w_2$  are $b^-$-vertices, which is not possible.

\noindent {\bf Case 2.} $w_1= z_1$ and $z_2=w_2$.

\noindent Since $\delta(D) = 2$, $N^+(z_1)= \{y_1, y_2\}$ and $N^-(x_1)= \{y_3, y_4\}$ where $N^+(z_1), N^-(x_1) \subseteq B\setminus \{x_2, x_4\}$. 

Let $\Gamma: V(D) \rightarrow \{c_1, c_2, c_3\}$ defined as: $\Gamma(x_4)  = c_1$; $\Gamma(x_1)  = \Gamma(x_2)  = \Gamma(z_1) = c_2$ and $\Gamma(x_3) = \Gamma(z_2) = c_3$.  For the vertices $\{y_1, y_2, y_3, y_4\}$, it is possible that either $\{y_1,y_2\}= \{y_3, y_4\}$; or $\{y_1,y_2\}\cap \{y_3, y_4\}\emptyset$ or $y_1=y_3$ and $y_2\not=y_4$. In any case, color the vertices with colors $c_1$ and $c_3$ in a way that both colors are present in $N^-(x_1)$ an $N^+(z_1)$. 

Also, for every $x\in A\setminus\{z_1, z_2, x_1, x_3, x_5\}$, $\Gamma(x) = c_2$ and for every $x\in B\setminus\{x_2, x_4, y_1, y_2, y_3, y_4\}$, $\Gamma(x) = c_1$. It is not hard to see that $\Gamma$ is a $b$-coloring of $D$ with 3 colors, where  $x_4$ is a  $b$-vertex;  $x_3$ and $z_1$ are $b^+$-vertices; and $x_1$ and $z_2$  are $b^-$-vertices, which is not possible. \end{proof}

From Claim \ref{b4} we see that if $dib(D)\leq 2$, then there is no bad path of order $4$. Since $\delta(D)\geq 2$, this is not possible, and the result follows. 

\end{proof}

\end{document}
